\subsection{Proofs for Section~\ref{sec:ext-game}}\label{app:proofs-game}

Notation is as in Section~\ref{sec:ext-game}. Throughout, $u(t):=(r+\delta_F)F_A(t)$ is the user cost, $I(a,b)=\int_a^be^{-(r-G)s}ds$, and $G+\delta_F>0$. Two facts are used repeatedly.

\begin{enumerate}
\item[(G1)] (Single crossing.) For $\mathcal E>0$, the sign of $\mathcal Ee^{Gt}-u(t)$ is the sign of $\mathcal Ee^{(G+\delta_F)t}-(r+\delta_F)F_0$, which is strictly increasing in $t$. Hence $\mathcal Ee^{Gt}<u(t)$ for $t<T^\circ(\mathcal E)$ and $\mathcal Ee^{Gt}>u(t)$ for $t>T^\circ(\mathcal E)$, and $T^\circ$ is nonincreasing in $\mathcal E$, strictly where it is positive. For $\mathcal E\le0$, $\mathcal Ee^{Gt}<u(t)$ at every date.
\item[(G2)] (Single-firm rule.) A firm whose adoption raises its flow by $\mathcal Ee^{Gt}$ from date $t_0$ on, whatever else happens, maximizes $\int_T^\infty e^{-rs}\mathcal Ee^{Gs}ds-e^{-rT}F_A(T)$ over $T\ge t_0$ at $\max\{t_0,T^\circ(\mathcal E)\}$. The derivative with respect to $T$ is $-e^{-rT}[\mathcal Ee^{GT}-u(T)]$, so this is Proposition~\ref{prop:ext-timing}(b) with $W(t)=\mathcal Ee^{Gt}$, together with (G1).
\end{enumerate}

\paragraph{Lemma~\ref{lem:ext-game-pass}.} (a) Write $\mathrm{mg}_0(k):=\mathrm{mg}-k\bar\varphi^\times\delta_c$ and $\mathrm{mg}_1(k):=\mathrm{mg}_0(k)+(1-\bar\varphi^{\rm own})\delta_c$. Both are positive for $k\le n_c-1$, because $\mathrm{mg}>(n_c-1)\bar\varphi^\times\delta_c$ and $\bar\varphi^{\rm own}<1$, and both are strictly decreasing in $k$ because $\bar\varphi^\times>0$. So $\mathcal E_a(k)=\mathrm{mg}_a(k)^2/B_q$ is strictly decreasing in $k$, and $\mathcal E^{\rm priv}(k)=(1-\bar\varphi^{\rm own})\delta_c[2\mathrm{mg}_0(k)+(1-\bar\varphi^{\rm own})\delta_c]/B_q$ is positive and strictly decreasing. Cournot: by Lemma~\ref{lem:ext-retained}(a), firm $i$'s own cut raises its margin by $n_c\delta_c/(n_c+1)$, each rival's cut lowers it by $\delta_c/(n_c+1)$, and profit is margin squared over $B$; this gives the stated values, with $\bar\varphi^{\rm ind}=1-n_c/(n_c+1)+(n_c-1)/(n_c+1)=n_c/(n_c+1)$. Bertrand: the first-order condition for $p_i$ is $q_i=p_i-c_i$, so profit is $(p_i-c_i)^2$ and $B_q=1$. The best responses are $(2+s_B)p_i=A_B+c_i+s_B\sum_lp_l$. Summing, $\sum_lp_l=(n_cA_B+\sum_lc_l)/[2-(n_c-1)s_B]$, which is well defined because $(n_c-1)s_B<1$. Hence
\[
p_i-c_i=\frac{A_B+s_B\sum_lp_l-(1+s_B)c_i}{2+s_B},\qquad \frac{\partial(p_i-c_i)}{\partial c_l}=\bar\varphi^\times\ (l\ne i),\qquad \frac{\partial(p_i-c_i)}{\partial c_i}=\bar\varphi^\times-\frac{1+s_B}{2+s_B},
\]
with $\bar\varphi^\times=s_B/\{(2+s_B)[2-(n_c-1)s_B]\}=s_B\bar\varphi^{\rm ind}/(2+s_B)>0$. Margins are therefore linear in the cuts, a firm's own cut raises its margin by $(1-\bar\varphi^{\rm own})\delta_c$ with $\bar\varphi^{\rm own}=1/(2+s_B)+\bar\varphi^\times$, and each rival's cut lowers it by $\bar\varphi^\times\delta_c$. A common change in costs moves every price by $1/[2-(n_c-1)s_B]$, which equals $1-(1-\bar\varphi^{\rm own})+(n_c-1)\bar\varphi^\times$.

(b) Firm $i$'s first-order condition is $P(Q)+P'(Q)q_i=c_i$. Summing over firms, $n_cP+P'Q=\sum_lc_l$. The derivative of the left side in $Q$ is $(n_c+1)P'+P''Q=P'(n_c+1-\mathsf e)$, which is nonzero at a regular equilibrium, so by the implicit function theorem the equilibrium is continuously differentiable in costs near the symmetric point. A common change $dc$ gives $dQ=n_c\,dc/[P'(n_c+1-\mathsf e)]$ and $dP=P'dQ$, so $\bar\varphi^{\rm ind}=n_c/(n_c+1-\mathsf e)$. For a change $dc_i$ in one firm's cost, the rivals' conditions give $P'dQ+P''Y\,dQ+P'dq_r=0$, that is $dq_r=-(1-\mathsf e/n_c)dQ$ for each rival, and the summed condition gives $dQ=dc_i/[P'(n_c+1-\mathsf e)]$. So $dP=dc_i/(n_c+1-\mathsf e)$, and $\bar\varphi^{\rm own}=\bar\varphi^{\rm ind}/n_c$. Firm $i$'s earnings are $\mathcal E_i=(P-c_i)q_i$, with $P-c_i=-P'Y$ at the symmetric point and $dq_i=dQ-(n_c-1)dq_r=[n_c-(n_c-1)\mathsf e/n_c]dQ$. Hence
\[
\frac{d\mathcal E_i}{dc_i}=Y\Bigl(\frac{\bar\varphi^{\rm ind}}{n_c}-1\Bigr)-Y\,\frac{n_c-(n_c-1)\mathsf e/n_c}{n_c+1-\mathsf e}=-Y\Bigl[\frac{2n_c-1}{n_c}-\frac{(n_c-1)\bar\varphi^{\rm ind}}{n_c^2}\Bigr],
\]
using $\mathsf e=n_c+1-n_c/\bar\varphi^{\rm ind}$. Because the equilibrium is continuously differentiable, $\mathcal E^{\rm priv}(k)$ equals $-\delta_c$ times this derivative at the symmetric point plus $o(\delta_c)$, for every $k$, since the $k$ rivals' cuts move the evaluation point only by $O(\delta_c)$. For a common change, each firm earns $(P-c)Q/n_c$, and $(P-c)\,dQ/n_c=-P'Y\bar\varphi^{\rm ind}dc/(n_cP')=-Y\bar\varphi^{\rm ind}dc/n_c$, so its derivative in $c$ is $Y(\bar\varphi^{\rm ind}-1)-Y\bar\varphi^{\rm ind}/n_c=-Y[1-\bar\varphi^{\rm ind}(1-1/n_c)]$, which gives $\mathcal E^{\rm eq}$. Summed over firms, this is the producer-surplus formula $dPS/dc=-Q[1-(1-\theta)\bar\varphi^{\rm ind}]$ of \citet{weyl2013pass} with $\theta=1/n_c$. Finally, $\varphi=(dP/d\hat c)(\hat c/P)$ by Lemma~\ref{lem:prices}(a).

(c) With $P=(Q/D)^{-1/\epsilon}$, $P'Q/P=-1/\epsilon$, so the condition $P+P'q_i=c_i$ reads $P(1-1/(n_c\epsilon))=\hat c$ at the symmetric point, which requires $n_c\epsilon>1$ and gives $P=n_c\epsilon\hat c/(n_c\epsilon-1)$. Then $\bar\varphi^{\rm ind}=n_c\epsilon/(n_c\epsilon-1)$, $\varphi=1$, and the Lerner index is $1/(n_c\epsilon)$. Revenue per firm is $DP^{1-\epsilon}/n_c$ and earnings are the Lerner index times revenue, so both are proportional to $\hat c^{\,1-\epsilon}$; replacing $\hat c$ by $\hat c-\delta_c$ multiplies them by $(1-\delta_c/\hat c)^{1-\epsilon}$, which is below one if and only if $\epsilon<1$. Here $\mathsf e=1+1/\epsilon$. Substituting $\bar\varphi^{\rm ind}$ into the bracket of the private gain in \eqref{eq:ext-game-first} and multiplying by $n_c^2(n_c\epsilon-1)>0$ shows that it is positive if and only if $(2n_c-1)(n_c\epsilon-1)>(n_c-1)\epsilon$, that is, $\epsilon(2n_c^2-2n_c+1)>2n_c-1$. \qed

\paragraph{Proposition~\ref{prop:ext-game-preempt}.} (a) The dates are nondecreasing because $\mathcal E^{\rm priv}$ is decreasing and $T^\circ$ nonincreasing (G1), and strictly increasing where positive. Fix the firm assigned $T_k$ and the others' dates. When $m$ of its rivals have adopted, its adoption raises its flow by $\mathcal E^{\rm priv}(m)e^{Gt}$, and the derivative of its payoff with respect to its own date $T$ is $-e^{-rT}[\mathcal E^{\rm priv}(m(T))e^{GT}-u(T)]$, where $m(T)$ counts the rivals that have adopted by $T$. For $T\in[T_j,T_{j+1})$ with $j<k$, $m(T)=j$, and $\mathcal E^{\rm priv}(j)e^{GT}<u(T)$ because $T<T_{j+1}=T^\circ(\mathcal E^{\rm priv}(j))$ (G1). For $T\in[T_k,T_{k+1})$, $m(T)=k-1$ and $T\ge T^\circ(\mathcal E^{\rm priv}(k-1))$. For $T\in[T_j,T_{j+1})$ with $j>k$, $m(T)=j-1$ and $T\ge T_j=T^\circ(\mathcal E^{\rm priv}(j-1))$. So the derivative is positive before $T_k$ and nonpositive after it, and $T_k$ is a best response. The last claim follows from (G1), since $T^\circ$ is strictly decreasing where it is positive.

(b) By Assumption~\ref{ass:ext-game}, $\mathcal E_0(1)<\mathcal E_0(0)$ and $\mathcal E_1(1)<\mathcal E_1(0)$, so $\mathcal E_1(0)-\mathcal E_0(1)>\mathcal E^{\rm priv}(0)>\mathcal E^{\rm priv}(1)$. For $t<T_2$, differentiating \eqref{eq:ext-game-values},
\[
\frac{d}{dt}\bigl[V^{\rm L}(t)-V^{\rm F}(t)\bigr]=-e^{-rt}\bigl[(\mathcal E_1(0)-\mathcal E_0(1))e^{Gt}-u(t)\bigr],
\]
which by (G1) is positive before $T_D$ and negative after it. At $T_2$ the two values coincide. If $T_D$ were zero, $V^{\rm L}-V^{\rm F}$ would fall on $[0,T_2]$ to zero and be positive at 0, against $V^{\rm L}(0)<V^{\rm F}(0)$; so $T_D>0$, and then $T_D<T_1<T_2$ by (G1). $V^{\rm L}-V^{\rm F}$ is negative at 0 and positive at $T_D$, so it has a unique root $T^{\rm P}$ in $(0,T_D)$ and no other root in $[0,T_2)$. Three further properties are used. (i) For $t<T_2$, $V^{\rm L}{}'(t)=-e^{-rt}[\mathcal E^{\rm priv}(0)e^{Gt}-u(t)]$, so $V^{\rm L}$ is strictly increasing on $[0,T_1]$, and $V^{\rm L}(T^{\rm P})<V^{\rm L}(T_1)$. (ii) For $t<T_2$, $V^{\rm F}(t)-V^{\rm J}(t)=\int_t^{T_2}e^{-rs}[u(s)-\mathcal E^{\rm priv}(1)e^{Gs}]ds>0$: simultaneous adoption is worse than following. (iii) $dV^{\rm F}/dt=e^{-(r-G)t}[\mathcal E_0(0)-\mathcal E_0(1)]>0$ on $[0,T_2)$. After a first adoption at $s$, the other firm's flow rises by $\mathcal E^{\rm priv}(1)e^{Gt}$ when it adopts, whatever happens later, so by (G2) it adopts at $\max\{s,T_2\}$; this justifies \eqref{eq:ext-game-values}. \citet{fudenberg1985preemption} construct, in their extended space of strategies for timing games (the limits of discrete-time mixed strategies), a subgame-perfect equilibrium for payoffs with this structure: the follower's best response just described, $V^{\rm L}<V^{\rm F}$ before $T^{\rm P}$, $V^{\rm L}>V^{\rm F}>V^{\rm J}$ on $(T^{\rm P},T_2)$, $V^{\rm L}$ increasing up to $T^{\rm P}$, and continuous payoffs. In it, neither firm adopts before $T^{\rm P}$, and in every subgame that starts at $t\ge T^{\rm P}$ with no adoption, the firms adopt with the intensities of their construction, which give each firm $V^{\rm F}(t)$. At $T^{\rm P}$, where $V^{\rm L}=V^{\rm F}$, one firm adopts and simultaneous adoption has probability zero. Their argument uses only these properties of the leader, follower and joint payoffs (their $L$, $F$ and $M$), so it applies to \eqref{eq:ext-game-values}. A deviation to adopt at $t<T^{\rm P}$ yields $V^{\rm L}(t)<V^{\rm L}(T^{\rm P})$ by (i), and waiting past $T^{\rm P}$ cannot yield more than $V^{\rm F}(T^{\rm P})$. On the equilibrium path, each firm therefore earns $V^{\rm L}(T^{\rm P})=V^{\rm F}(T^{\rm P})$. The verification script also solves the discrete-time game by backward induction with symmetric stage equilibria, and the first adoption converges to $T^{\rm P}$ as the period shrinks (check B6). \qed

\paragraph{Proposition~\ref{prop:ext-game-wait}.} (a) $\mathcal E^{\rm eq}=\mathcal E_1(1)-\mathcal E_0(0)<\mathcal E_1(1)-\mathcal E_0(1)=\mathcal E^{\rm priv}(1)$, so $T_J>T_2$ by (G1), and $T_J=\infty$ if $\mathcal E^{\rm eq}\le0$. $V^{\rm J}{}'(t)=-e^{-rt}[\mathcal E^{\rm eq}e^{Gt}-u(t)]$, so $V^{\rm J}$ is strictly increasing before $T_J$ and strictly decreasing after it. Consider the strategies: adopt at $T$ if no one has adopted; after a first adoption at $s$, adopt at $\max\{s,T_2\}$. The response is optimal by (G2). A deviation to adopt first at $t<T_2$ yields $V^{\rm L}(t)\le\max_{[0,T_2]}V^{\rm L}=V^{\rm L}(T_1)$, by the derivative in the previous proof and (G1). A deviation to adopt first at $t\in[T_2,T)$ yields $V^{\rm J}(t)\le V^{\rm J}(T)$, because the rival then adopts at once. Delaying past $T$ changes nothing, since the rival adopts at $T$ and the deviator then follows at once. So the strategies form a subgame-perfect equilibrium if and only if $V^{\rm J}(T)\ge V^{\rm L}(T_1)$; the subgames off the path are single-firm problems. Joint adoption at $T<T_2$ is not an equilibrium, because waiting to follow yields $V^{\rm F}(T)>V^{\rm J}(T)$ (property (ii) above). Joint adoption at finite $T>T_J$ is not one either, because adopting at any $t\in(T_J,T)$ yields $V^{\rm J}(t)>V^{\rm J}(T)$. Such equilibria exist if and only if $\sup_{T\in[T_2,T_J]}V^{\rm J}(T)=V^{\rm J}(T_J)\ge V^{\rm L}(T_1)$, with $V^{\rm J}(\infty)=\mathcal E_0(0)/(r-G)$. Since $V^{\rm J}(T_2)=V^{\rm L}(T_2)$, this is $V^{\rm J}(T_J)-V^{\rm J}(T_2)\ge V^{\rm L}(T_1)-V^{\rm L}(T_2)$, and the two sides are the integrals of the derivatives in \eqref{eq:ext-game-joint}. Finally, $V^{\rm J}(T)\ge V^{\rm L}(T_1)>V^{\rm L}(T^{\rm P})$.

(b) In the good state the payoffs are those of \eqref{eq:ext-game-values}, with the follower adopting at $\max\{t,T_2\}$. In the bad state the leader keeps $\mathcal E_0(0)$ and loses its outlay, and the follower keeps $\mathcal E_0(0)$. This gives $V^{\rm L}_\varpi$ and $V^{\rm F}_\varpi$. For $t\ge T_2$, $V^{\rm L}_\varpi-V^{\rm F}_\varpi=-(1-\varpi)e^{-rt}F_A(t)<0$. For $t<T_2$,
\[
V^{\rm L}_\varpi(t)-V^{\rm F}_\varpi(t)=\varpi\bigl[\mathcal E_1(0)-\mathcal E_0(1)\bigr]I(t,T_2)+\varpi e^{-rT_2}F_A(T_2)-e^{-rt}F_A(t),
\]
whose derivative is $-e^{-rt}[\varpi(\mathcal E_1(0)-\mathcal E_0(1))e^{Gt}-u(t)]$. If $\varpi\le\bar\varpi$, then $\varpi(\mathcal E_1(0)-\mathcal E_0(1))\le\mathcal E^{\rm priv}(1)$, so the bracket is negative on $[0,T_2)$ by (G1). The difference is then increasing on $[0,T_2]$ up to $-(1-\varpi)e^{-rT_2}F_A(T_2)<0$. Also $\bar\varpi<1$ because $\mathcal E_1(0)>\mathcal E_1(1)$.

(c) For $t<T_2$, $V^{\rm L}_\varpi{}'(t)=-e^{-rt}[\varpi\mathcal E^{\rm priv}(0)e^{Gt}-u(t)]$. Here $\varpi\mathcal E^{\rm priv}(0)<\varpi(\mathcal E_1(0)-\mathcal E_0(1))\le\mathcal E^{\rm priv}(1)$, so the derivative is positive by (G1). For $t>T_2$, $V^{\rm L}_\varpi(t)=\varpi V^{\rm J}(t)+(1-\varpi)[\mathcal E_0(0)/(r-G)-e^{-rt}F_A(t)]$, with derivative $-e^{-rt}[\varpi\mathcal E^{\rm eq}e^{Gt}-u(t)]$, positive before $T^{\rm W}$ and negative after it. Since $\varpi\mathcal E^{\rm eq}<\mathcal E^{\rm priv}(1)$, $T^{\rm W}\ge T_2$, and $V^{\rm L}_\varpi$ is continuous at $T_2$. (i) Let the designated follower never adopt first and, after a first adoption at $s$, adopt at $\max\{s,T_2\}=s$ if the bundle works and never otherwise; this response is optimal by (G2). The leader's best date is then the maximizer $T^{\rm W}$ of $V^{\rm L}_\varpi$. Never adopting is worse, because $V^{\rm L}_\varpi(T^{\rm W})-\mathcal E_0(0)/(r-G)=e^{-rT^{\rm W}}[\varpi e^{GT^{\rm W}}\mathcal E^{\rm eq}/(r-G)-F_A(T^{\rm W})]$. This is positive: if $T^{\rm W}>0$, $\varpi\mathcal E^{\rm eq}e^{GT^{\rm W}}=u(T^{\rm W})$ and $(r+\delta_F)/(r-G)>1$; if $T^{\rm W}=0$, $\varpi\mathcal E^{\rm eq}\ge(r+\delta_F)F_0$. The follower cannot gain by leading at $t<T^{\rm W}$, which yields $V^{\rm L}_\varpi(t)<V^{\rm L}_\varpi(T^{\rm W})<V^{\rm F}_\varpi(T^{\rm W})$, nor by adopting together with the leader, which adds the outlay in the bad state. Off the path, if the designated leader has not adopted by $t>T^{\rm W}$, adopting at once is its best choice, because $V^{\rm L}_\varpi$ is decreasing there. (ii) Let the rival adopt first with hazard $\nu^{\rm W}$ and distribution function $\Gamma$. For $s\ge T^{\rm W}$, a firm that plans to adopt at $s$ if the rival has not adopted earns $U(s)=\int_{T^{\rm W}}^sV^{\rm F}_\varpi\,d\Gamma+(1-\Gamma(s))V^{\rm L}_\varpi(s)$. Then $U'(s)=(1-\Gamma(s))[\nu^{\rm W}(s)(V^{\rm F}_\varpi-V^{\rm L}_\varpi)+V^{\rm L}_\varpi{}']=0$, because $V^{\rm F}_\varpi-V^{\rm L}_\varpi=(1-\varpi)e^{-rs}F_A(s)$ and $V^{\rm L}_\varpi{}'=-e^{-rs}[\varpi\mathcal E^{\rm eq}e^{Gs}-u(s)]$, which gives \eqref{eq:ext-game-hazard}. So $U(s)=U(T^{\rm W})=V^{\rm L}_\varpi(T^{\rm W})$ for all $s\ge T^{\rm W}$. A plan to adopt at $s<T^{\rm W}$ yields $V^{\rm L}_\varpi(s)<V^{\rm L}_\varpi(T^{\rm W})$, because the rival does not adopt before $T^{\rm W}$. Every date in the support is therefore a best response, and the mixed strategy is an equilibrium; the same hazard applies in every later subgame. The hazard is nonnegative after $T^{\rm W}$ and grows without bound, because $e^{(G+\delta_F)t}\to\infty$, so $\Gamma\to1$. The first of the two independent adoption times exceeds $T^{\rm W}+u$ with probability $\exp\{-2\int_{T^{\rm W}}^{T^{\rm W}+u}\nu^{\rm W}\}$. With $T^{\rm W}>0$, $\int_{T^{\rm W}}^{T^{\rm W}+u}\nu^{\rm W}=(r+\delta_F)[(e^{(G+\delta_F)u}-1)/(G+\delta_F)-u]/(1-\varpi)$. This is the symmetric equilibrium of a war of attrition \citep{hendricks1988war}. \qed

\paragraph{Proposition~\ref{prop:ext-game-pd}.} (a) Adopting is strictly dominant if and only if $\mathcal E_1(k)-\mathcal K>\mathcal E_0(k)$ whatever the number $k$ of rival adopters, that is, $\mathcal K<\min_k\mathcal E^{\rm priv}(k)=\mathcal E^{\rm priv}(n_c-1)$. When all adopt, each firm earns $\mathcal E_1(n_c-1)-\mathcal K$, which is at most $\mathcal E_0(0)$ if and only if $\mathcal K\ge\mathcal E^{\rm eq}$. The width is $\mathcal E^{\rm priv}(n_c-1)-\mathcal E^{\rm eq}=\mathcal E_0(0)-\mathcal E_0(n_c-1)>0$. Dividing \eqref{eq:ext-game-first} by $Y\delta_c$ gives \eqref{eq:ext-game-window}, and with $\bar\varphi^{\rm ind}=n_c/(n_c+1)$ its ends are $2/(n_c+1)$ and $2n_c/(n_c+1)$, whose ratio is $1/n_c$.

(b) By \eqref{eq:ext-game-first}, the first-order common gain is negative if and only if $\bar\varphi^{\rm ind}>n_c/(n_c-1)$. With $\bar\varphi^{\rm ind}=n_c/(n_c+1-\mathsf e)$ and $n_c+1-\mathsf e>0$, this is $n_c-1>n_c+1-\mathsf e$, that is, $\mathsf e>2$, whatever $n_c$. Then $\mathcal E^{\rm eq}<0\le\mathcal K$ for every $\mathcal K\ge0$, and (a) gives the dilemma whenever $\mathcal K<\mathcal E^{\rm priv}(n_c-1)$. Under constant elasticity, Lemma~\ref{lem:ext-game-pass}(c) gives the exact statement. Under linear demand $\mathsf e=0$.

(c) Before $T_{n_c}$ the last adopter earns $\mathcal E_0(m)\le\mathcal E_0(0)$ when $m$ rivals have adopted, and from $T_{n_c}$ on it earns $\mathcal E_1(n_c-1)$. Its gain over never adopting is therefore at most $e^{-(r-G)T_{n_c}}\mathcal E^{\rm eq}/(r-G)-e^{-rT_{n_c}}F_A(T_{n_c})$, which is nonpositive if and only if $\mathcal E^{\rm eq}\le(r-G)e^{-GT_{n_c}}F_A(T_{n_c})$. If $T_{n_c}>0$, then $\mathcal E^{\rm priv}(n_c-1)e^{GT_{n_c}}=u(T_{n_c})$, which gives \eqref{eq:ext-game-dyn}. In the duopoly each firm earns $V^{\rm F}(T^{\rm P})<V^{\rm F}(T_2)=V^{\rm J}(T_2)$, by property (iii) in the proof of Proposition~\ref{prop:ext-game-preempt}. Moreover, $V^{\rm J}(T_2)-\mathcal E_0(0)/(r-G)$ is the same expression with $T_2$ in place of $T_{n_c}$ and $\mathcal E^{\rm priv}(1)$ in place of $\mathcal E^{\rm priv}(n_c-1)$. Under linear demand, $\mathcal E^{\rm eq}/\mathcal E^{\rm priv}(n_c-1)\to1/n_c$ as $\delta_c\to0$ by (a), so to first order \eqref{eq:ext-game-dyn} is $1/n_c\le(r-G)/(r+\delta_F)$.

(d) With $m_0=\mathcal E_0(0)/S_0$ and $m_1=(\mathcal E_1(n_c-1)-\mathcal K)/S_1$, $m_1-m_0=[\mathcal E_1(n_c-1)-\mathcal K-\mathcal E_0(0)S_1/S_0]/S_1$, which is the stated identity. Revenue per firm is $PQ/n_c$, so $d\ln S=(1-\epsilon)\,d\ln P$ and $d\ln P=\bar\varphi^{\rm ind}\,d\hat c/P=\varphi\,d\hat c/\hat c$, with $d\hat c=-\delta_c$. The margin before $\mathcal K$ is the gross margin $1-\hat c/P$, and Lemma~\ref{lem:prices}(a) gives its change $(1-\varphi)\delta_c/P$. A continuing fall in unit cost at rate $\delta_u$ makes the price fall at $\varphi\delta_u$, and Lemma~\ref{lem:prices}(b) gives the change in nominal growth. Under constant elasticity, earnings and revenue are multiplied by the same factor (Lemma~\ref{lem:ext-game-pass}(c)), so $\mathcal E_1(n_c-1)/S_1=m_0$ and $m_1=m_0-\mathcal K/S_1$. The adoption cost raises the cost floor $a$ by $\mathcal K/S_1$, so the margin intercept $m_0=\gamma-a$ of Part~I falls by that amount, and $v(s)$ cannot rise, because the payoff of every spend policy in \eqref{eq:value} falls. Nominal market size, which is industry revenue, is multiplied by $(1-\delta_c/\hat c)^{1-\epsilon}<1$ when $\epsilon<1$, and $V=Mv(s)$ by Proposition~\ref{prop:dynamic}(a). \qed

\paragraph{Proposition~\ref{prop:ext-game-entry}.} (a) $P_k\ge P_{n_c}$ for every $k$, and $P_E$ is decreasing, so for $t>\bar T_M$, $P_k>P_E(t)$ whatever the number of adopters, and entry binds. An incumbent then earns $(P_E(t)-\hat c_0)Ye^{Gt}$ before adopting and $(P_E(t)-\hat c_0+\delta_c)Ye^{Gt}$ after, whatever its rivals do. The gain $\delta_cYe^{Gt}$ does not depend on prices, so (G2) gives the date. (b) $P_k$ is decreasing in $k$ and $P_E$ is strictly decreasing, so $T_M(k)$, the date after which $P_E<P_k$, is nondecreasing in $k$. With $k(t)$ adoptions by $t$, entry never binds if and only if $P_{k(t)}\le P_E(t)$ at every $t$. Let $k^*(t):=\min\{k:P_k\le P_E(t)\}$, which is at most $\bar k$ if $P_{n_c}\le P_E^\infty$. Since $P$ is decreasing in $k$, the condition is $k(t)\ge k^*(t)$ for all $t$. Now $k^*(t)\ge j$ if and only if $P_{j-1}>P_E(t)$, that is, $t>T_M(j-1)$. So the condition holds if and only if $k(t)\ge j$ for every $t>T_M(j-1)$ and every $j\le\bar k$, which, since $k(\cdot)$ is right-continuous, means that the $j$-th adoption takes place by $T_M(j-1)$. (c) In the long run, $P_E(t)\to P_E^\infty$. An adopter earns $(P_E^\infty-\hat c_0+\delta_c)Y-\mathcal K$ a year, scaled by $e^{Gt}$, a non-adopter $(P_E^\infty-\hat c_0)Y$, and before AI the firm earned $(P_0-\hat c_0)Y=\mathcal E_0(0)$. The two differences are as stated. Adoption is a best response but leaves the firm no better off than before AI, the defensive case of Proposition~\ref{prop:ext-should}(a), exactly when the first difference is positive and the second nonpositive. If $P_E^\infty\le P_0-\delta_c$, the lower end of the window is nonpositive, so every $\mathcal K\in[0,\delta_cY)$ lies in it. The viability statements compare long-run unit margins with zero. (d) The long-run unit margin of an adopter is $P_E^\infty-(\hat c_0-\delta_c)=p_{\rm sw}-[(\hat c_0-\delta_c)-\hat c_E^\infty]\le p_{\rm sw}$. A scarce input lowers the incumbent's cost relative to entrants', as in Proposition~\ref{prop:entry}(c). \qed

\paragraph{Corollary~\ref{cor:ext-game-engine}.} (a) Set $n_c=1/\mathrm{HHI}$ in \eqref{eq:ext-game-first}. The common share is $1-\bar\varphi^{\rm ind}(1-\mathrm{HHI})$, and with $\bar\varphi^{\rm ind}=n_c/(n_c+1-\mathsf e)$ it equals $(2-\mathsf e)/(n_c+1-\mathsf e)$. The private share is $(2n_c-1)/n_c-(n_c-1)\bar\varphi^{\rm ind}/n_c^2=2-\mathrm{HHI}-\mathrm{HHI}(1-\mathrm{HHI})\bar\varphi^{\rm ind}$. Linear demand has $\bar\varphi^{\rm ind}=1/(1+\mathrm{HHI})$, and constant elasticity has $\bar\varphi^{\rm ind}=\epsilon/(\epsilon-\mathrm{HHI})$, which give the stated forms. If $n_r$ rivals replicate, the firm's gain is its earnings at the cost vector with its own and $n_r$ rivals' cuts, less $\mathcal E_0(0)$. By the differentiability in the proof of Lemma~\ref{lem:ext-game-pass}(b), to first order this is the own effect, $\mathrm{priv}\cdot Y\delta_c$, plus $n_r$ times the effect of one rival's cut. Setting $n_r=n_c-1$ identifies that effect as $-(\mathrm{priv}-\mathrm{dur})Y\delta_c/(n_c-1)$, and $\mathrm{rep}=n_r/(n_c-1)$ gives the interpolation. With constant $\mathsf e$, $d\,\mathrm{dur}/d\mathrm{HHI}=(2-\mathsf e)/[\mathrm{HHI}^2(1/\mathrm{HHI}+1-\mathsf e)^2]$, which has the sign of $2-\mathsf e$, and $\mathrm{dur}=1$ at $\mathrm{HHI}=1$. Where entry binds, Proposition~\ref{prop:ext-game-entry}(c) gives the gains over the status quo and over not adopting, as shares of $\delta_cY$. (b) The definitions apply Proposition~\ref{prop:ext-game-pd}(a) and Definition~\ref{def:ext-motives} to the risk-adjusted present values, with the gain from replicable savings measured after the rivals' response. (c) This is (a) multiplied by $\varpi\delta_cY$, less $\mathcal K$; as $\mathrm{HHI}\to0$ with constant $\mathsf e$, $\mathrm{dur}\to0$. \qed
